\section{Introdução}

Este trabalho trata da previsão de renda: dada uma pessoa, descrita por 10 atributos do censo dos EUA de 1994, prever se a sua renda anual passa de 50 mil dólares. A base Adult~\cite{uci:adult,kohavi:96} tem cerca de 30 mil pessoas depois da curadoria. Nenhum atributo separa as classes sozinho, há atributos categóricos e valores ausentes, e o tamanho da base já torna visível o custo de treinar uma máquina de vetores de suporte (SVM)~\cite{boser:92,cortes:95} com kernel.

O modelo principal é a SVM, em quatro versões (linear, com kernel RBF, com kernel polinomial e RBF com pesos por classe), comparada com cinco modelos de referência (a classe majoritária, uma regra de um corte, a regressão logística, a árvore de decisão e a Random Forest) e com uma rede neural. Além da acurácia, registramos a AUC, os erros em cada classe e o tempo de treino, de busca de hiperparâmetros e de previsão de cada modelo. As contribuições são:
\begin{itemize}
  \item uma curadoria que mede o erro inevitável da base: perfis idênticos com rendas diferentes limitam qualquer modelo a \svTeto\% de acurácia (Seção~\ref{sec:sv_dados});
  \item um estudo da SVM numa divisão: kernels, o mapa de $C$ e $\gamma$, o efeito da padronização e o crescimento do tempo de treino com o tamanho da base (Seção~\ref{sec:sv_estudo});
  \item uma comparação em 10 seeds, com acurácia, F1, AUC, erros por classe e tempo num único núcleo, que mostra que a busca guiada pela acurácia piora a AUC da SVM com kernel RBF (Seção~\ref{sec:sv_seeds}).
\end{itemize}
